\documentclass{jsarticle}
\usepackage[dvipdfmx,bookmarks=true]{hyperref}
\usepackage{pxjahyper}
\usepackage{pdfpages}

\hypersetup{
 setpagesize=false,
 bookmarksnumbered=true,%
 colorlinks=true,%
 linkcolor=blue,
 citecolor=red,
}
\usepackage{ascmac}

\title{不明点に対する回答書}
\author{\empty}
\date{\empty}

\begin{document}
\maketitle
\noindent \textbf{論文番号：Q38} \\\
\noindent \textbf{標題：コストと精度を考慮した質問と文書の特徴に基づく適応的なRAG構成選択}  \\


\newpage

\section*{査読者A様の査読結果への回答}

\subsection*{A-1}
\begin{screen}
近接する適応的RAG／routing系研究との比較実験が存在しない。関連研究はSelf-RAG、AutoRAG-HP、AdaGCRAG等に触れるにとどまり、第3節で挙げた近接研究との比較がない。比較対象が固定構成・ランダム・オラクルに限られているため、本手法が既存の適応的RAGより優れる、あるいは同等精度で低コストであるという含意を支持できない。少なくとも代表的な1手法(質問複雑度で分岐するAdaptive-RAG [Jeong 24] は最も比較が容易である)との実測比較を追加するか、比較が困難である理由を技術的に説明すべきである。
\end{screen}

\subsection*{A-1への回答}
改訂論文の2節および表1において，近接研究（CARROT，uRAG，RAGRouter，AdaGCRAG）
と本研究との関係性を整理し，提案手法のポジショニングを明確化しました．
4.5節，4.6節，5.2節，5.4節，表4，表5において，Adaptive-RAG [Jeong 24] との実測比較を追加しました．

\subsection*{A-2}
\begin{screen}
固定ベースラインに対する優位性が示されていない。表3において、最頻出設定はF1=43.55／1.41秒であるのに対し、提案@1のF1値-時間はF1=43.58／1.43秒であり、実質的に同等である。さらに提案@1の時間重視は最頻出設定と完全に一致している。したがって、抄録および結論の「固定構成に対しF1改善・時間短縮を両立」という主張は、LangChain標準設定との比較に限ってのみ成立し、より強い固定ベースラインに対しては成立していない。本文でこの事実を率直に記述し、主張を適切に限定すべきである。加えて、本研究の貢献を立証するには、最頻出設定および既存適応手法を上回るPareto改善(同等コストでの精度向上、または同等精度での低コスト化)を実証する必要がある。
\end{screen}

\subsection*{A-2への回答}
改訂論文の5.2節および表5において，全クエリに同一のRAG 構成を適用する固定構成（標準設定・最頻出設定）や，既存の代表的な適応的RAG 手法であるAdaptive-RAG と比較して，回答精度（EM，Token-level F1 値）の大幅な改善（Pareto 改善）があることを加筆しました．

\subsection*{A-3}
\begin{screen}
評価項目(iii)のコスト計上が論文の主題と整合しない。4.1節では「ユーザが上位K個の候補を実行し、最良を採用する」運用を想定しているが、表4で報告される時間は採用された1構成の時間のみであり、K個を試行した総コストが計上されていない。本論文はコスト感度を主要な貢献に掲げているため、運用を想定するなら累積実行時間や累積推論コストを併記すべきである。現状では、ユーザ選択による精度向上が追加コストに見合うかを判断できない。
\end{screen}

\subsection*{A-3への回答}
改訂論文の5.3節および表6において，上位 K 件提示モデルにおける構成順次試行コストを「累積実行時間」として計上し，より実用的なコスト効率性を実証しました．

\subsection*{A-4}
\begin{screen}
文書特徴の寄与が限定的であり、主要主張と整合しない。表5では、RAG構成(C)単独でNDCG@1=66.84であるのに対し、全特徴を用いてもNDCG@1=68.52にとどまり、改善幅は約1.7ポイントである。さらにC+Q+S(64.94)はC単独を下回るなど、多くの組合せで追加特徴が性能向上に寄与していない。この結果からは「文書側の特徴を統合することで高精度なランキングが可能になる」という貢献の第一点を強く支持できない。どの条件下で文書特徴が有効に働くのかを分析し、改善幅の限定性についても正直に議論すべきである。
\end{screen}

\subsection*{A-4への回答}
表8（旧表5）のアブレーション分析結果に対するご指摘を受け，RAG構成情報（C）がランキング予測における支配的な基盤要因であること，および文書・質問特徴量の寄与と改善幅（約1.7ポイント）の限定性について，論文全体の主張と分析を全面的に見直し，以下の改訂を行いました．

\begin{enumerate}
\item \textbf{貢献の主張の適正化（第1節，第6節）：}
第1節の主要な貢献（箇条書き第1点）および第6節の記述について，「文書側の特徴を統合することで高精度なランキングが可能になる」という表現を改め，生成モデルや検索方式といった「RAG構成情報自身が順位付けの基盤要因となる一方で，質問・文書の特徴量を包括的に統合することで構成単独（NDCG@1=66.84\%）を上回る最高精度（NDCG@1=68.52\%，Hit@5=65.00\%）を達成できる」という，実験事実に基づいた限定的かつ正確な主張に修正しました．

\item \textbf{アブレーション分析の深化と条件下での有効性の分析（第5.5節）：}
5.5節において，以下の分析・考察を加筆しました．
\begin{itemize}
\item \textbf{C単独の支配的役割：} RAG構成（C）単独がNDCG@1=66.84\%，Hit@5=62.50\%，NDCG@5=68.33\%（全条件中最良）を達成しており，マクロな構成要素（生成モデルや検索設定等）が候補構成の性能の上界・下界を大枠で規定していることを明記しました．
\item \textbf{部分的な特徴量追加における性能低下の要因：} Cに対して一部の特徴量のみを追加した組み合わせ（例: C+Q: 63.81\%，C+D: 60.66\%，C+Q+S: 64.94\%等）がC単独を下回る現象について，質問の意図・構文構造や文書特性との相互作用を捉えきれない部分的な特徴量付与が決定木（LightGBM）においてノイズとして働き，過学習や不適切な分岐を生じさせるためであることを説明しました．
\item \textbf{文書特徴が有効に機能する条件：} 質問の複雑さ（質問タイプや構文木の深さ）と，文書の情報量（トークン数）および意味的一致度（コサイン類似度・固有表現類似度）が同時に提供されることで，例えば「構文が複雑で文書トークン数が大きくエンティティ重複が高い状況ではGraphRAGや大チャンク構成を選択し，単純な事実抽出質問に対しては局所検索や軽量構成を選択する」といった精緻な条件分岐が可能となり，全特徴量を統合した場合（C+Q+D+S）に最高精度（NDCG@1=68.52\%，Hit@5=65.00\%）が達成されることを示しました．
\item \textbf{改善幅の限定性の率直な議論：} 全特徴量を統合した場合の改善幅が約1.7ポイントにとどまり，NDCG@5ではC単独とほぼ同等である点を正直に記述し，データ側の特徴量は最良構成（Top-1）の絞り込みを微細に補正する補助的役割を担っていることを明確化しました．
\end{itemize}

\item \textbf{文書特徴量の表現能力に関する限界の明記（第5.8節）：}
現在の文書特徴量がトークン数や大域的埋め込み類似度に依存しており，文書内の局所的な論理構造や複数文書間のトポロジーを直接捉えられていないことが，改善幅が約1.7ポイントにとどまる要因であると限界事項に明記し，今後の課題としてグラフ構造特徴やチャンク間情報密度分布などの高次文書特徴量の抽出・統合の必要性を加筆しました．
\end{enumerate}
\begin{screen}
統計的有意性の検定がない。評価データは80事例と小規模であり、報告される性能差は1〜3ポイント程度である。にもかかわらず、本文5.8節等で「有意な差」という表現が用いられているが、信頼区間や有意差検定は示されていない。ブートストラップ法または対応のある検定により、報告された改善が統計的に意味を持つかを確認すべきである。
\end{screen}

\subsection*{A-5への回答}
改訂論文の4.6節および5.4節において，提案手法とベースライン間の差について，Wilcoxon 符号付き順位検定による有意差検定について加筆し，回答精度の統計的優位性を立証しました．

\subsection*{A-6}
\begin{screen}
精度指標がBERTScoreのみである。SQuADやHotpotQA等には標準的なEM／token-F1が存在するため、これらを併記すべきである。BERTScoreは意味的類似性の評価に優れる一方で誤答に部分点を与える傾向があり、定性分析で示される誤答例(例: 正解mount holyoke collegeに対しtemple university)が過大評価されている可能性を排除できない。標準QA指標との傾向の一致を確認することで、結論の信頼性が高まる。
\end{screen}

\subsection*{A-6への回答}
改訂論文の4.6節，5.2節，表5において，自動評価指標（BERTScore）に加えて，標準的QA評価指標であるExact Match (EM)およびToken-level $F_1$値を用いて提案手法を多角的に再評価しました．

\subsection*{A-7}
\begin{screen}
総合評価値(式1)の集約方式に対する検証がない。5.4節は逆数ランクの単純加算が有効と結論しているが、線形和や対数等の代替集約との比較がなく、定式化の妥当性が示されていない。簡単なアブレーションの追加が望ましい。
\end{screen}

\subsection*{A-7への回答}
改訂論文の5.5節および表7において，式(1) で示した総合評価値の集約方式の妥当性を評価するため，(1) 逆数ランク和（提案手法），(2) 対数順位和（Log-Rank），(3) 予測値線形和（Linear Score）を用いたアブレーション実験について加筆しました．本実験結果より，特別なスケーリングを必要とせず，上位候補の識別能力に優れる逆数ランク和を用いた式(1) の単純加算の定式化は，精度とコストの柔軟な制御と実運用上の簡便性を両立するうえで最も有効な設計であると結論付けられました．

\subsection*{A-8}
\begin{screen}
重み設定と挙動の不整合。表4の提案@3では、F1値重視(再現率59.58)が再現率重視(59.37)より再現率が高い。重み設計の直感に反する現象であり、説明を要する。また表4の上位K候補がどの重み設定で生成されたかが4.1節からは追えないため、明記すべきである。
\end{screen}

\subsection*{A-8への回答}
表6（旧表4）における上位$K$候補の抽出条件の不明瞭さ，および$F_1$値重視と再現率重視における数値挙動について，以下の通り論文を修正・加筆いたしました．

\begin{enumerate}
\item \textbf{上位$K$候補の抽出条件の明記（第4.0節，第4.3節）：}
評価項目(iii)における実験手順について，各重視方針（適合率重視，再現率重視，$F_1$値重視，時間重視）に対応する重み設定（$W_{\mathrm{pre}}=1.0$，$W_{\mathrm{rec}}=1.0$，$W_{\mathrm{f_1}}=1.0$，$W_{\mathrm{time}}=1.0$）を用いてRMにより各質問の上位$K$個の候補構成を抽出し，ユーザがそれら$K$個を実行した結果から当該方針に対応する指標が最良となる1構成を採用する運用フローであることを第4.0節に明確に追記しました．

\item \textbf{提案@3における再現率の挙動に関する説明の加筆（第5.3節）：}
提案@3において$F_1$値重視の再現率（59.58\%）が再現率重視（59.37\%）をわずかに（0.21ポイント）上回っている要因について，以下の通り5.3節に考察を加筆しました．
\begin{itemize}
\item 各重視方針に対応する重み設定でRMが事前予測した上位$K$候補の中から事後選択を行うため，RMの予測誤差に伴い抽出される$K$個の構成プールが方針ごとに異なること．
\item $F_1$値重視のRMは調和平均の最大化を目指して情報量の多い長めのコンテキストを生成する構成（大チャンク構成等）を上位に含めやすく，マルチホップ質問等の複雑な事例において正解トークンの網羅性（再現率）が副次的に高まりやすいためであること．
\end{itemize}
なお，両方針とも適合率重視（56.67\%）や時間重視（55.18\%）と比較して再現率は大幅に向上しており，方針に応じた大枠のトレードオフ制御は意図通りに機能していることを確認しております．
\end{enumerate}
\subsection*{A-9}
\begin{screen}
質問タイプ判定が英語固有のルールベースである。is/can/must、what/who/where、orといった英語の語に依存しており、評価データも全て英語である。本誌読者が関心を持つ日本語や専門ドメインでの一般化可能性について、5.8節の限界に言語依存性を明記すべきである。
\end{screen}

\subsection*{A-9への回答}
改訂論文の5.9節に言語依存性を明記しました．

\textbf{第五に，質問タイプ分類は is, can, what, who, or などの英語の疑問詞・助動詞に依存するルールベースであり，日本語の質問にそのまま適用することはできない．
日本語環境へ拡張するには，MeCab などを用いた形態素解析により疑問詞や係助詞を判定し，言語ごとに質問タイプの定義とルールを再設計する必要がある．
また，専門ドメインでは一般語の疑問詞だけでは意図を十分に捉えられない場合があるため，ドメイン固有の質問分類器への置換も今後の課題である．}

\subsection*{A-10}
\begin{screen}
ランダム選択の数値が異常である。表3のランダム選択は適合率0.06\%、再現率23.37\%と乖離が極端で、多数の構成が失敗している可能性がある。算出方法(質問ごとの平均か全体平均か)を含め、原因を説明すべきである。
\end{screen}

\subsection*{A-10への回答}
表4（旧表3）におけるランダム選択の数値表記および算出方法・乖離原因について，以下の通り修正および加筆を行いました．

\begin{enumerate}
\item \textbf{表記誤りの修正（第5.2節，表4）：}
旧原稿において，ランダム選択の適合率を小数値（0.065）からパーセント表記へ換算する際の転記誤りにより「0.06\%」と誤って記載しておりました．正しくは全7,680試行の全体平均値である「適合率: 6.54\%，再現率: 24.12\%，$F_1$値: 8.92\%，時間: 50.95秒」であり，表4および本文5.2節の数値を正確な値に修正いたしました．

\item \textbf{算出方法の明確化（第4.5節）：}
ランダム選択のベースラインの算出方法について，80問のテスト質問に対して全96通りの構成空間から一様無作為に構成を選択した際の期待値（全7,680試行の全体平均）であることを第4.5節に明記しました．

\item \textbf{適合率が著しく低く再現率と乖離する原因の技術的説明（第5.2節）：}
ランダム選択において適合率（6.54\%）が再現率（24.12\%）に比べて大幅に低迷し，全試行の81.6\%で適合率が0となる原因について，以下の通り5.2節に考察を加筆しました．
\begin{itemize}
\item \textbf{構成空間の偏りと失敗応答：} 全96構成のうち3分の2（64通り）を占めるGraphRAGやLightRAG，Map-Reduce等の重厚な構成は，単純な事実抽出質問に対しても長大な推論過程や定型的な謝絶文（例:「提供データからは判断できません」等）を出力する傾向が強く，正解との形式的不一致やハルシネーションにより適合率が0となる構成が多数存在すること．
\item \textbf{BERTScoreの非対称性：} 長大な推論テキスト中に正解トークンが一部含まれている場合，再現率（Recall）はある程度獲得できるものの，膨大な不要トークンに対する強い長さペナルティにより適合率（Precision）が極端に低下すること．
\end{itemize}
この結果は，構成を無作為に適用すると深刻な構成ミスマッチと大幅な遅延（平均50秒超）が生じることを示しており，質問や文書の特徴に応じて最適なRAG構成を適応的に選択する提案手法の必要性を強く裏付けるものであると議論を整理しました．
\end{enumerate}
\subsection*{A-11}
\begin{screen}
会議版からの差分の明示。本論文は言語処理学会版 [穴口 26] の改訂版であるため、技術論文として加えられた増分(ablation・特徴量・定性分析が該当するのであればその旨)を冒頭で具体的に記述すべきである。
\end{screen}

\subsection*{A-11への回答}

「1. はじめに」の最後に，下記の記述を追加しました．

\textbf{なお，本論文は，穴口らの研究成果 [穴口 26] に対し，以下の具体的な追加・拡張を行ったものである．
第一に，近接研究（CARROT，uRAG，RAGRouter，AdaGCRAG）との関係性を整理し，手法のポジショニングを明確化した．
第二に，既存の代表的な適応的RAGベースライン手法であるAdaptive-RAGを評価フレームワークに組み込み，精度とコストのParetoフロンティアの観点から性能評価を行った．
第三に，自動評価指標（BERTScore）に加えて，標準的QA評価指標であるExact Match (EM)およびToken-level $F_1$値を用いて提案手法を多角的に再評価した．
第四に，提案手法とベースライン間の差について，Wilcoxon符号付き順位検定による有意差検定を行い，回答精度の統計的優位性を立証した．
最後に，上位K件提示モデルにおける構成順次試行コストを「累積実行時間」として計上し，より実用的なコスト効率性を実証した．}

\subsection*{A-12}
\begin{screen}
英文抄録の校正。「base don」「con-figuration」「consistingof」「demonstrated thatselecting」「toall」「ablation,feature」など改行由来の語結合・誤記が散見される。本文5.3節の数値(再現率59.46)も表4(59.58)と不一致である。
\end{screen}

\subsection*{A-12への回答}
改訂論文における英語抄録では，改行由来の語結合・誤記等を修正いたしました．

\newpage

\section*{査読者B様の不明点への回答}

\subsection*{B-1: 妥当性}
\begin{screen}
提案方法の新規性としては，式（１）及び式（２）と予想される．
しかし，式（２）の効果が実験結果から適正に評価されているかが不明である．
各重みを線形的に結合することの妥当性が不明である．
例を挙げると，表３で適合率を重視しているものが，
他を重視しているものよりも適合率が低いことの理由が不明である．

単に，その項目を重視してランキングの高い組合せを優先するのであれば，
式（１）の各項だけを用いると良いように思う．
実際に結合を評価しているのは，第一の設定のみである．
結合の妥当性を評価するのであれば，複数の重みに変化をつけ，
その結果がその優先度の高い項が高くなるように重みを最適化する必要があるように考える．

結論として，この実験方法及び実験結果からは式（１）と式（２）の妥当性が不明である．
\end{screen}

\subsection*{B-1への回答}
ご指摘ありがとうございます．

式(1)および式(2)の重み付き線形結合の妥当性評価，および表4（旧表3）における適合率重視設定の挙動について，以下の通り論文の加筆・修正を行いました．

\begin{enumerate}
\item \textbf{重み結合の妥当性評価（第5.4.1節，表7）：}
査読者様のご指摘の通り，複数の重みを連続的に変化させて結合効果を評価するため，改訂論文の第5.4.1節および表7において，重み $W_{\mathrm{time}}$ を $0.0$ から $1.0$ まで $0.1$ 刻みで変化させた「重みスイープ分析」を追加しました．
実験結果より，$W_{\mathrm{time}}$ の増加に伴い実行時間が 7.24秒から 1.41秒へと段階的かつ単調に短縮され，それに伴い $F_1$値が 45.08\% から 43.55\% へと滑らかに遷移するという，明確な多目的トレードオフ制御が機能することを実証しました．
さらに，代替集約方式として (1) 逆数ランク和（提案手法），(2) 対数順位和（Log-Rank），(3) 予測値線形和（Linear Score）を比較し，事前のスケーリングを要さず上位候補の識別能力に優れる式(1)・式(2)の定式化が，精度とコストの柔軟な制御と実運用上の簡便性を両立するうえで最も有効であることを確認しました．

\item \textbf{適合率重視設定における適合率の挙動と多目的統合の必要性（第5.2節）：}
表4（旧表3）において適合率重視設定（$W_{\mathrm{pre}}=1.0$）の実測適合率（39.16\%）が他設定を下回る要因について，第5.2節に以下の考察を加筆しました．
\begin{itemize}
\item 適合率（Precision）は生成回答中の余分なトークンに対して極めて敏感であり，データセット全体の81.6\%で適合率が0となるなど学習ラベルがスパースでノイズが大きいこと．
\item 適合率単独（$W_{\mathrm{pre}}=1.0$）のRMは文脈脱落を過度に警戒して大きめのチャンクサイズ（CS=500）を選定しやすく，結果として生成文が長文化して適合率を落とす局所解に陥ること．
\end{itemize}
この現象は，まさに単一指標（単項）のみに依存した最適化の難しさを示しており，適合率と再現率の調和平均である$F_1$値や実行時間といった複数指標を統合する式(1)・式(2)の多目的定式化が，局所解を回避する正則化として不可欠であることを示しています．実際，表7の重みスイープに示す通り，$F_1$値と時間を組み合わせた複合設定では安定した高精度・低遅延の最適化が達成されております．
\end{enumerate}
\subsection*{B-2: 有用性}
\begin{screen}
著者らは，５．８節で限界を示している．
第二の部分の問題で，「適合率やF1値の最大化」について，
オラクル構成と有意な差がある状態の結果は，
実際の利用において有用と言えるのかが不明である．
\end{screen}

\subsection*{B-2への回答}
オラクル構成との差が存在する状況における提案手法の実用的な「有用性」について，以下の通り論文（第5.8節，第5.4節）を改訂し，考察を明確化いたしました．

\begin{enumerate}
\item \textbf{オラクルの位置付けの明確化（第5.8節）：}
本研究で定義した「オラクル」は，各質問に対して全96通りのRAG構成を実際にすべて実行した結果から事後的に最良値を選択する「理論上の上限値（全探索オラクル）」です．実運用において全探索を行うことは，1クエリあたり約50秒以上の遅延および96倍のAPI計算コストを要するため現実的ではなく，あくまでモデル予測性能の上限を示す参照基準として設定しております．

\item \textbf{実用ベースラインに対する統計的優位性と有用性の立証（第5.2節，第5.4節）：}
実運用において開発者が採用可能な現実的ベースラインとの比較において，提案手法（提案@1）は以下の明確な優位性（有用性）を達成しています．
\begin{itemize}
\item \textbf{固定構成に対する優位性：} 一般に用いられるLangChain標準設定（$F_1$: 40.50\%，時間: 5.18秒）と比較して，$F_1$値を 45.08\% へと大幅に向上させつつ時間を 1.43秒へと大幅に短縮（72\%の時間削減）しています．さらに最強の固定構成である最頻出設定に対しても，EM（30.00\% vs 27.50\%）およびToken-level $F_1$（45.51\% vs 39.25\%）において明確な改善を示しています．
\item \textbf{既存の適応手法（Adaptive-RAG）に対する優位性：} 代表的な適応的RAGベースラインであるAdaptive-RAG（$F_1$: 29.43\%，時間: 10.33秒）に対し，提案手法は$F_1$値で約14--16ポイント上回り，かつ時間は約7分の1（1.43秒）に短縮しています．Wilcoxon符号付き順位検定により，この精度改善は統計的に極めて有意（$p < 0.001$）であることが立証されています．
\item \textbf{推論オーバーヘッドの極小性：} LLMによるルーティング判断を推論時に介在させるエージェント型手法とは異なり，提案手法は決定木によるミリ秒オーダーの特徴量推論のみで構成を決定するため，実運用における応答遅延や追加コストがほぼゼロです．
\item \textbf{上位$K$件提示によるオラクルへの漸近（第5.3節）：} 上位$K$件（$K=3, 5$）をユーザに提示して事後選択を行う運用モデルでは，$F_1$値が 50.53\% / 50.25\%，実行時間が 1.27秒（オラクルの時間最良設定と完全一致）に到達し，オラクルとの差を実用的に大幅に縮小できることを示しています．
\end{itemize}
以上より，全探索オラクルとの差は単一推論によるランキング予測の今後の発展課題として存在しつつも，実運用可能なすべてのベースライン手法に対して統計的有意なPareto改善を達成していることから，本手法の実用的な有用性は十分に確立されていると結論付けております．
\end{enumerate}

\end{document}
